% ==============================================================================
% lab1_siem_detection.tex — UE SIS589
% Travaux Pratiques 1 : Déploiement SIEM, Corrélation et Chasse aux Alertes
% ==============================================================================

\chapter{TP~1~— Déploiement SIEM, Corrélation et Chasse aux Alertes}
\label{lab:siem}

\epigraph{\itshape « Un SIEM sans règles calibrées est un radar sans opérateur. »}
         {--- Principe SOC}

% ── Présentation ──────────────────────────────────────────────────────────────
\section*{Présentation du TP}

\begin{table}[H]
\centering\small
\rowcolors{2}{TableauPair}{TableauImpair}
\begin{tabular}{p{3.5cm}p{10cm}}
\toprule
\textbf{Durée} & 4~heures (2~h déploiement + 2~h analyse) \\
\textbf{Objectifs} & Déployer Wazuh + ELK~; configurer les agents~; créer
  des règles de corrélation~; qualifier des alertes réelles~; produire
  un rapport d'alerte structuré \\
\textbf{Pré-requis} & Administration Linux niveau opérationnel~;
  notions de réseau TCP/IP~; chapitres~1--3 lus \\
\textbf{Infrastructure} & Environnement GNS3/EVE-NG fourni par l'enseignant.
  Accès SSH à srv-siem-01 (Wazuh Manager)~; agents pré-installés sur
  srv-web-01, srv-erp-01, dc-01 \\
\textbf{Livrable} & Rapport de TP~: captures d'écran des tableaux de bord,
  règles rédigées, 3~alertes qualifiées, recommandations \\
\bottomrule
\end{tabular}
\end{table}


% ==============================================================================
\section{Partie~A — Vérification et prise en main de l'infrastructure SIEM}
\label{lab1:partieA}
% ==============================================================================

\subsection{A1 — Connexion et état de l'architecture}

Connectez-vous au serveur SIEM de LiZbiZ-SIS~:

\begin{lstlisting}[style=StyleTerminal,
  caption={Connexion initiale au SIEM et vérification des services},
  label={lst:lab1-connexion}]
# Connexion SSH au serveur Wazuh Manager
ssh etudiant@10.0.0.5  # Mot de passe fourni en séance

# ─── Vérifier que tous les services sont actifs ───────────────────────────────
sudo systemctl status wazuh-manager
sudo systemctl status elasticsearch
sudo systemctl status kibana
sudo systemctl status logstash

# ─── Lister les agents Wazuh connectés ───────────────────────────────────────
sudo /var/ossec/bin/agent_control -l
# Résultat attendu :
# ID: 001, Name: srv-web-01,   IP: 10.0.1.10, Active
# ID: 002, Name: srv-erp-01,   IP: 10.0.2.10, Active
# ID: 003, Name: dc-01,         IP: 10.0.3.10, Active

# ─── Vérifier la réception de logs en temps réel ─────────────────────────────
sudo tail -f /var/ossec/logs/alerts/alerts.log | head -20
\end{lstlisting}

\begin{boxexo}
\textbf{Question A1.1~:} Quelle est la différence entre un agent Wazuh
en état \texttt{Active} et \texttt{Disconnected}~? Quels événements
peuvent provoquer le passage d'un agent en état \texttt{Disconnected}~?

\medskip\textbf{Question A1.2~:} Consultez le fichier
\texttt{/var/ossec/etc/ossec.conf} sur le manager. Identifiez~:
(a)~la fréquence de vérification FIM~; (b)~les répertoires surveillés
sur les agents Linux~; (c)~le niveau minimal d'alerte transmis à Elasticsearch.
\end{boxexo}

\subsection{A2 — Navigation dans Kibana}

Ouvrez Kibana dans un navigateur~: \texttt{http://10.0.0.5:5601}

\begin{enumerate}
  \item Naviguez vers \textbf{Security~→~Overview}.
        Notez le nombre d'alertes des dernières 24~heures, le niveau
        moyen de criticité, et les top~3 agents les plus actifs.
  \item Ouvrez \textbf{Discover}. Sélectionnez l'index
        \texttt{wazuh-alerts-*}. Filtrez sur les 12~dernières heures.
  \item Identifiez le champ \texttt{rule.level} et filtrez pour afficher
        uniquement les alertes de niveau~$\geq 8$.
  \item Exportez les colonnes \texttt{@timestamp},
        \texttt{agent.name}, \texttt{rule.id}, \texttt{rule.description},
        \texttt{data.srcip} en CSV (bouton Share~→~CSV Export).
\end{enumerate}

\begin{boxexo}
\textbf{Question A2.1~:} Combien d'alertes de niveau~$\geq 10$ ont été
générées sur les 12~dernières heures~? Sur quel agent~? Quelle règle~?

\medskip\textbf{Question A2.2~:} Wazuh utilise une échelle de niveaux
de 0 à 15. À partir de quel niveau une alerte est-elle considérée
comme haute~? Citez la documentation Wazuh correspondante.
\end{boxexo}


% ==============================================================================
\section{Partie~B — Configuration des sources et des agents}
\label{lab1:partieB}
% ==============================================================================

\subsection{B1 — Activer la collecte Sysmon sur srv-web-01 (Windows)}

\begin{lstlisting}[style=StyleTerminal,
  caption={Déploiement et configuration Sysmon sur srv-web-01},
  label={lst:lab1-sysmon}]
# Sur srv-web-01 (session RDP ou PSExec depuis la salle TP)

# Étape 1 : Vérifier que Sysmon est installé
sc query Sysmon64
# Si absent : installer (fichiers fournis sur le partage réseau TP)
# \\10.0.0.1\TP-LAB1\Sysmon64.exe -accepteula -i sysmonconfig.xml

# Étape 2 : Charger la configuration sysmonconfig.xml fournie
# (configuration minimale du chapitre 2, adaptée LiZbiZ-SIS)
Sysmon64.exe -c \\10.0.0.1\TP-LAB1\sysmonconfig-lab1.xml

# Étape 3 : Vérifier que les événements Sysmon sont bien collectés
# par l'agent Wazuh (vérifier ossec.conf sur l'agent)
type "C:\Program Files (x86)\ossec-agent\ossec.conf" | findstr -i "sysmon"
# Doit contenir : <location>Microsoft-Windows-Sysmon/Operational</location>

# Étape 4 : Forcer l'envoi d'un événement test
# Ouvrir PowerShell et exécuter une commande encodée (simulée)
powershell -EncodedCommand ZQBjAGgAbwAgAHQAZQBzAHQA
# (= "echo test" en base64 — génère un Sysmon EID 1)
\end{lstlisting}

\begin{boxexo}
\textbf{Question B1.1~:} Dans Kibana, confirmez que l'événement Sysmon
EID~1 généré par la commande PowerShell est bien visible dans l'index
Wazuh. Copiez l'entrée JSON complète dans votre rapport.

\medskip\textbf{Question B1.2~:} Identifiez dans le JSON les champs~:
\texttt{process.name}, \texttt{process.command\_line}, \texttt{process.hash.sha256},
\texttt{process.parent.name}. Quelle valeur contient \texttt{rule.mitre.id}~?
\end{boxexo}

\subsection{B2 — Activer la collecte Apache sur srv-web-01 (Linux)}

\begin{lstlisting}[style=StyleTerminal,
  caption={Configuration de la collecte des logs Apache via Wazuh Agent},
  label={lst:lab1-apache}]
# Sur srv-web-01 Linux (SSH)
# Vérifier la configuration de l'agent Wazuh
sudo cat /var/ossec/etc/ossec.conf | grep -A5 "localfile"

# Ajouter la collecte Apache si absente
sudo tee -a /var/ossec/etc/ossec.conf << 'EOF'
  <localfile>
    <log_format>apache</log_format>
    <location>/var/log/apache2/access.log</location>
  </localfile>
  <localfile>
    <log_format>apache</log_format>
    <location>/var/log/apache2/error.log</location>
  </localfile>
EOF

# Redémarrer l'agent
sudo systemctl restart wazuh-agent

# Générer du trafic test vers le serveur web
curl -s http://10.0.1.10/ > /dev/null
curl -s "http://10.0.1.10/?id=1 UNION SELECT 1,2,3--" > /dev/null
curl -s "http://10.0.1.10/wp-login.php" -X POST \
  -d "log=admin&pwd=test123" > /dev/null
for i in $(seq 1 30); do
  curl -s "http://10.0.1.10/$(cat /dev/urandom | tr -dc 'a-z' | \
    head -c8).php" > /dev/null
done
# Les 30 dernières requêtes génèrent des 404 en masse

# Vérifier dans Kibana que les alertes Apache apparaissent
\end{lstlisting}

\begin{boxexo}
\textbf{Question B2.1~:} Dans Kibana, filtrez les alertes provenant de
srv-web-01 sur les 15~dernières minutes. Identifiez les règles Wazuh
déclenchées. Quel est le Rule ID associé à la détection de SQLi~?

\medskip\textbf{Question B2.2~:} La requête avec \texttt{UNION SELECT}
a-t-elle déclenché une alerte de niveau~$\geq 8$~? Si non, expliquez
pourquoi et proposez une modification de la configuration Wazuh pour
améliorer cette détection.
\end{boxexo}


% ==============================================================================
\section{Partie~C — Rédaction et test de règles de corrélation}
\label{lab1:partieC}
% ==============================================================================

\subsection{C1 — Règle de détection du brute force SSH}

\begin{lstlisting}[style=StyleCode, language=YAML,
  caption={Règle Wazuh à créer --- brute force SSH personnalisée},
  label={lst:lab1-regle-ssh}]
<!-- Fichier : /var/ossec/etc/rules/local_rules.xml -->
<!-- À compléter et tester -->

<group name="lab1_lizbiz,">

  <!-- Règle 1 : Brute force SSH (5 échecs en 60s, même IP) -->
  <rule id="100501" level="10" frequency="5" timeframe="60">
    <if_matched_sid>5716</if_matched_sid>
    <!-- TODO : ajouter le champ de corrélation par IP source -->
    <!-- TODO : ajouter la description -->
    <!-- TODO : ajouter le mapping MITRE T1110.001 -->
    <!-- TODO : ajouter le groupe pci_dss_10.2.4 -->
  </rule>

  <!-- Règle 2 : Authentification SSH réussie après des échecs -->
  <!-- (connexion réussie dont le même sid 5716 a été déclenché avant) -->
  <rule id="100502" level="12">
    <!-- TODO : définir la règle parent (brute force suivi de succès) -->
    <!-- TODO : mapper MITRE T1078 (Valid Accounts) -->
  </rule>

</group>
\end{lstlisting}

\begin{boxexo}
\textbf{Travail~C1~:}
\begin{enumerate}
  \item Complétez les deux règles XML ci-dessus (remplacez tous
        les \texttt{TODO}).
  \item Chargez les règles~: \texttt{sudo /var/ossec/bin/wazuh-logtest -V}
        pour valider la syntaxe.
  \item Testez la règle~100501 en générant un brute force SSH depuis
        la VM attaquante~: \\
        \texttt{hydra -l root -P /usr/share/wordlists/rockyou.txt \\ ssh://10.0.1.10 -t 4 -f}
  \item Confirmez l'alerte dans Kibana. Copiez le JSON de l'alerte.
  \item Calculez le FPR de votre règle sur les 30~dernières minutes.
\end{enumerate}
\end{boxexo}

\subsection{C2 — Règle de détection de mouvement latéral SMB}

\begin{lstlisting}[style=StyleCode, language=YAML,
  caption={Règle à créer --- mouvement latéral SMB (T1021.002)},
  label={lst:lab1-regle-smb}]
<!-- Mouvement latéral : connexion SMB réussie depuis un poste vers
     plusieurs serveurs en moins de 5 minutes (T1021.002) -->
<rule id="100510" level="11" frequency="3" timeframe="300">
  <!-- TODO : identifier les Event IDs Windows pertinents (EID 4624) -->
  <!-- TODO : filtrer sur LogonType 3 (réseau) -->
  <!-- TODO : regrouper par compte source (not by IP) -->
  <!-- TODO : exclure les comptes de service légitimes -->
  <!-- TODO : mapper T1021.002 -->
</rule>
\end{lstlisting}

\begin{boxexo}
\textbf{Travail~C2~:} Complétez, chargez et testez la règle~100510.
Simulez un mouvement latéral depuis dc-01 vers srv-erp-01 et srv-web-01
avec les credentials fournis. Documentez~: (a)~l'alerte générée~;
(b)~le délai entre l'action et la détection (MTTD sur cet événement)~;
(c)~une amélioration possible.
\end{boxexo}


% ==============================================================================
\section{Partie~D — Qualification d'alertes et tableau de bord}
\label{lab1:partieD}
% ==============================================================================

\subsection{D1 — Qualification de 3~alertes}

À partir des alertes générées dans les parties précédentes, qualifiez
3~alertes selon le modèle suivant~:

\begin{table}[H]
\centering\small
\caption{Modèle de qualification d'alerte — à remplir pour chaque alerte}
\label{tab:lab1-qualification}
\begin{tabular}{p{4cm}p{9.5cm}}
\toprule
\textbf{Champ} & \textbf{Valeur} \\
\midrule
Identifiant alerte & \textit{(ID ticket)} \\
Timestamp (UTC) & \textit{(copier depuis le JSON)} \\
Agent source & \textit{(nom de l'agent Wazuh)} \\
Rule ID / Description & \textit{(ID + texte)} \\
IP source & \textit{(si disponible)} \\
Utilisateur & \textit{(si disponible)} \\
Niveau / Gravité & \textit{(niveau Wazuh + gravité IR)} \\
TTP MITRE ATT\&CK & \textit{(tactique + technique)} \\
Verdict & ☐ Vrai positif \quad ☐ Faux positif \quad ☐ À analyser \\
Justification du verdict & \textit{(2--3 phrases)} \\
Action recommandée & \textit{(confinement / surveillance / fermeture)} \\
\bottomrule
\end{tabular}
\end{table}

\subsection{D2 — Construction du tableau de bord SOC}

Dans Kibana~→~Dashboard~→~Create new dashboard, créez un tableau de
bord \texttt{SOC-LiZbiZ-SIS} comprenant les visualisations suivantes~:

\begin{enumerate}
  \item \textbf{Compteur~:} nombre total d'alertes de niveau $\geq 8$
        sur les dernières 24~heures.
  \item \textbf{Camembert~:} répartition des alertes par niveau
        (6~à~15).
  \item \textbf{Histogramme temporel~:} volume d'alertes par heure sur
        48~heures (pour détecter les pics d'activité).
  \item \textbf{Table Top~10~:} top~10 des adresses IP sources
        d'alertes, avec nombre d'alertes et niveau maximum.
  \item \textbf{Carte géographique~:} géolocalisation des IPs sources
        (si GeoIP configuré dans Logstash).
\end{enumerate}

\begin{boxexo}
\textbf{Question D2.1~:} À partir de votre tableau de bord, identifiez
l'heure à laquelle l'activité malveillante a débuté lors du TP.
Quelle visualisation vous a permis de la détecter le plus rapidement~?

\medskip\textbf{Question D2.2~:} Proposez une 6\textsuperscript{e}
visualisation qui aurait une valeur opérationnelle pour l'analyste N1
de LiZbiZ-SIS. Justifiez en 3~phrases.
\end{boxexo}


% ==============================================================================
\section{Partie~E — Chasse aux menaces (Threat Hunting)}
\label{lab1:partieE}
% ==============================================================================

\subsection{E1 — Hypothèse et investigation}

\textbf{Hypothèse~:} Des outils de reconnaissance Active Directory
ont été exécutés sur le réseau LiZbiZ-SIS (T1087, T1069, T1018).

\begin{lstlisting}[style=StyleTerminal,
  caption={Requêtes KQL de threat hunting --- recon AD},
  label={lst:lab1-hunting}]
# ─── Recherche 1 : exécutions de net.exe / net1.exe (recon AD) ───────────────
event.code: "1"
  AND (process.name: "net.exe" OR process.name: "net1.exe")
  AND process.command_line: (*"user" OR *"group" OR *"localgroup" OR *"view*)

# ─── Recherche 2 : requêtes LDAP depuis des postes non-serveurs ───────────────
event.code: "3"
  AND destination.port: 389
  AND NOT source.ip: "10.0.3.10"  # exclure DC-01 lui-même

# ─── Recherche 3 : PowerShell avec imports de modules AD suspects ─────────────
event.code: "4104"
  AND winlog.event_data.ScriptBlockText: (
    *"Get-ADUser" OR *"Get-ADGroup" OR *"Get-ADDomain" OR
    *"net user" OR *"BloodHound" OR *"SharpHound"
  )

# ─── Recherche 4 : Enum des partages réseau depuis un poste ──────────────────
event.code: "5140"  # Accès à un partage réseau
  AND NOT winlog.event_data.SubjectUserName: (*"$" OR *"SYSTEM")
\end{lstlisting}

\begin{boxexo}
\textbf{Travail~E1~:}
\begin{enumerate}
  \item Exécutez les 4~requêtes dans Kibana. Pour chacune~:
        notez le nombre de résultats et identifiez les événements
        les plus suspects.
  \item La 3\textsuperscript{e} requête a-t-elle des résultats~?
        Si oui, ouvrez l'entrée la plus récente et documentez~:
        IP machine, utilisateur, contenu du ScriptBlock.
  \item Sur la base de vos résultats, formulez une conclusion~:
        l'hypothèse initiale est-elle confirmée~? Quel TTP
        MITRE correspond~?
  \item Transformez la recherche la plus pertinente en règle Wazuh
        de niveau~10. Rédigez le XML complet.
\end{enumerate}
\end{boxexo}


% ==============================================================================
\section{Livrable et critères d'évaluation}
\label{lab1:livrable}
% ==============================================================================

\textbf{Rapport à remettre dans les 48~heures} (format PDF, 15~pages
maximum) comprenant~:

\begin{enumerate}
  \item \textbf{Introduction~:} description de l'infrastructure et de
        l'état initial du SIEM.
  \item \textbf{Configuration~:} fichiers de configuration Wazuh
        et règles complètes (XML), avec captures d'écran de validation.
  \item \textbf{Qualifications d'alertes~:} 3~fiches remplies selon
        le tableau~\ref{tab:lab1-qualification}.
  \item \textbf{Tableau de bord~:} capture d'écran commentée avec
        justification de chaque visualisation.
  \item \textbf{Threat Hunting~:} résultats des 4~requêtes, conclusion,
        règle créée.
  \item \textbf{Conclusion~:} 3~recommandations d'amélioration pour le
        SOC de LiZbiZ-SIS, métriques calculées (MTTD observé, FPR
        estimé des règles rédigées).
\end{enumerate}

\begin{table}[H]
\centering\small
\caption{Barème Lab~1}
\label{tab:lab1-bareme}
\rowcolors{2}{TableauPair}{TableauImpair}
\begin{tabular}{p{6cm}p{8cm}p{1.5cm}}
\toprule
\tableauHeader{Critère} & \tableauHeader{Indicateurs} &
\tableauHeader{Points} \\
\midrule
Configuration SIEM opérationnelle & Agents connectés, sources actives & 2 \\
Règles Wazuh correctes & Syntaxe XML valide, mapping MITRE, test positif & 3 \\
Qualification des alertes & Verdict justifié, TTP correct, action cohérente & 2 \\
Tableau de bord & 5~visualisations présentes, commentées & 2 \\
Threat Hunting & Requêtes exécutées, résultats documentés, règle créée & 1 \\
\bottomrule
\multicolumn{2}{r}{\textbf{Total}} & \textbf{10} \\
\end{tabular}
\end{table}
